\documentclass{report}
\usepackage{amsmath,amssymb}
\setlength{\parindent}{0mm}
\setlength{\parskip}{1em}
\begin{document}
\begin{center}
\rule{6in}{1pt} \
{\large Armando Coco \\
{\bf Second Order Multigrid Methods for Elliptic Problems with Discontinuous Coefficients on an Arbitrary Interface}}

University of Catania,\\
Department of Mathematics and Computer Science Viale Andrea Doria, 6 95125,\\
Catania, Italy
\\
{\tt coco@dmi.unict.it}\\
Giovanni Russo\end{center}

Elliptic equations with discontinuous coefficients across a
codimensional-one interface arise in several applications.
Interface-fitted grid methods are difficult to use in the case of moving
interfaces, because of the computationally expensive meshing procedures
needed at each time step. In such cases an approach treating the
interface as embedded in a Cartesian grid may be preferred.

Multigrid technique is one of the most efficient strategy to solve a
class of partial differential equations, using a hierarchy of
discretizations.
Several multigrid approaches exist in literature to treat the jumping
coefficient problem in 2D when the interface is aligned with the
Cartesian grid. We mention the method based on operator-dependent
interpolation (Brandt et al., Dendy et al., ...), where the interpolation
is carried out by exploiting the continuity of the flux instead of the
gradient of the solution, and the method based on Galerkin Coarse Grid
Operator (Stuben et al.).

In this work, we propose a multigrid technique to solve the Poisson
equation with discontinuous coefficients along an arbitrary interface,
described by a level-set function in a Cartesian grid.
A proper treatment of the transfer operators and the relaxation on the
boundary is proposed, in order to achieve the best convergence factor
(i.e.\ the convergence factor predicted by the Local Fourier Analysis for
the inner relaxation and in a rectangular domain). The convergence factor
does not depend on the size of the grid, nor on the jump in the
coefficient. Therefore, this solver can be employed for instance in a
fluid dynamic context, where a two-phase flow is modeled by the
incompressible Navier-Stokes equations and the jump in the coefficient
consists of the jump in the viscosity or the density between the two
fluids.

The method is second order accurate in the solution and the gradient.


\end{document}
